\subsection{Krull–Akizuki 定理}

引理1. 设 A 是诺特整环，其中每个非零素理想都是极大理想，M 是有限生成的挠 A-模。则 M 的长度 long \(_{A}\) (M) 是有限的。

由于 M 是挠模，与 M 相关的每个素理想都是 \(\neq(0)\)，因而都是极大理想。引理由第IV章第2节第5号命题7得到。

引理2. 设 A 为环，T 为 A-模，\((\mathrm{T}_{\iota})\) 为 T 的子模的右定向族，其并为 T。则 \(\operatorname{long}_{\mathbf{A}}(\mathbf{T}) = \sup (\operatorname{long}_{\mathbf{A}}(\mathbf{T}_{\iota}))\)。

\(\text{long}_{\mathbf{A}}(\mathbf{T}_{\iota}) \leqslant \text{long}_{\mathbf{A}}(\mathbf{T})\) 对所有 \(\iota\) 成立。若不存在整数大于 \(\text{long}_{\mathbf{A}}(\mathbf{T}_{\iota})\)，则引理显然成立，此时等式两边都是无穷大。否则，令 \(\iota_{0}\) 为使 \(\text{long}_{\mathbf{A}}(\mathbf{T}_{\iota})\) 取得最大值的指标；于是 \(\mathbf{T}_{\iota_{0}} = \mathbf{T}\)，因为族 \((\mathbf{T}_{\iota})\) 是定向的；于是本情形下断言得证。

评注。该证明不假设 A 可交换。

引理3. 设 A 是诺特整环，其中每个非零素理想都是极大理想，M 是秩为 r 的无挠 A-模，a 是 A 的非零元素。则 A/Aa 是有限长度的 A-模，并且：

\[
\operatorname{long} _ {\mathbf {A}} (\mathrm{M} / a \mathrm{M}) \leqslant r. \operatorname{long} _ {\mathbf {A}} (\mathrm{A} / \mathrm{A} a).\tag{6}
\]

引理1表明 \(\mathrm{long}_{\mathbf{A}}(\mathbf{A} / \mathbf{A}a)\) 是有限的。先在 \(\mathbf{M}\) 有限生成的情形证明 (6)。由于 \(\mathbf{M}\) 无挠且秩为 \(r\)，存在一个子模 \(\mathbf{L}\)，它属于 \(\mathbf{M}\) 并同构于 \(\mathbf{A}'\)，且 \(\mathbf{Q} = \mathbf{M} / \mathbf{L}\) 是有限生成的挠 A-模，因而具有有限长度（引理1）。对于每个整数 \(n \geqslant 1\)，典范满射 \(\mathbf{M} / a^n\mathbf{M} \to \mathbf{Q} / a^n\mathbf{Q}\) 的核等于 \((\mathbf{L} + a^n\mathbf{M}) / a^n\mathbf{M}\)，并同构于 \(\mathbf{L} / (a^n\mathbf{M} \cap \mathbf{L})\)；因为

\[
a ^ {n} \mathbf {L} \subset a ^ {n} \mathbf {M} \cap \mathbf {L},
\]

因此

\[
\begin{array}{r l} \operatorname{long} _ {\mathbf {A}} (\mathrm{M} / a ^ {n} \mathrm{M}) & \leqslant \\ \operatorname{long} _ {\mathbf {A}} (\mathrm{L} / a ^ {n} \mathrm{L}) + \operatorname{long} _ {\mathbf {A}} (\mathrm{Q} / a ^ {n} \mathrm{Q}) & \leqslant \operatorname{long} _ {\mathbf {A}} (\mathrm{L} / a ^ {n} \mathrm{L}) + \operatorname{long} _ {\mathbf {A}} (\mathrm{Q}). \end{array} \tag {7}
\]

现在，由于 M 无挠，乘以 a 定义了一个同构

将 M/aM 映到 \(aA/a^{2}M\)；对 L 也同样如此；由此对 n 归纳得到公式：

\[
\operatorname{long} _ {\mathbf {A}} (\mathbf {M} / a ^ {n} \mathbf {M}) = n. \operatorname{long} _ {\mathbf {A}} (\mathbf {M} / a \mathbf {M}).\tag{8}
\]

\[
\operatorname{long} _ {\mathbf {A}} (\mathbf {L} / a ^ {n} \mathbf {L}) = n. \operatorname{long} _ {\mathbf {A}} (\mathbf {L} / a \mathbf {L}).
\]

结合 (7)，得到：

\[
\operatorname{long} _ {\mathbf {A}} (\mathrm{M} / a \mathrm{M}) \leqslant \operatorname{long} _ {\mathbf {A}} (\mathrm{L} / a \mathrm{L}) + n ^ {- 1} \operatorname{long} _ {\mathbf {A}} (\mathrm{Q})\tag{9}
\]

对所有 \(n > 0\) 成立；由于 \(L\) 同构于 \(A^r\)，有 \(\text{long}_A(L/aL) = r \text{ long}_A(A/Aa)\)；令 (9) 中的 \(n\) 趋于无穷大，即得 (6)。

现在转到一般情形。设 \((\mathbf{M}_{\iota})\) 为 M 的有限生成子模族。模 T = M/aM 是子模 \(\mathrm{T}_{\iota} = (\mathrm{M}_{\iota} + a\mathrm{M})/a\mathrm{M} = \mathrm{M}_{\iota}/(\mathrm{M}_{\iota} \cap a\mathrm{M})\) 的并。现在，\(T_{\iota}\) 同构于 \(M_{\iota}/aM_{\iota}\) 的一个商，因而

\[
\operatorname{long} _ {\mathbf {A}} \left(\mathrm{T} _ {t}\right) \leqslant r \operatorname{long} _ {\mathbf {A}} (\mathrm{A} / \mathrm{A} a)
\]

由我们刚才证明的结果，因而

\[
\operatorname{long} _ {\mathbf {A}} (\mathrm{T}) \leqslant r \operatorname{long} _ {\mathbf {A}} (\mathrm{A/Aa})
\]

由引理2得到。

命题5（Krull–Akizuki）. 设 A 是诺特整环，其每个非零素理想都是极大理想，K 为其分式域，L 为 K 的有限扩张，B 是包含 A 的 L 的子环。则 B 是诺特环，且 B 的每个非零素理想都是极大理想。此外，对于 B 的每个理想 b ≠ (0)，B/b 是有限生成的 A-模。

设 b 是 B 的非零理想。我们将证明 B/b 是有限长度的 A-模（因而当然也是有限长度的 B-模），并且 b 是有限生成的 B-模。

b 的一个非零元素 y 满足如下形式的方程：

\[
a _ {r} y ^ {r} + a _ {r - 1} y ^ {r - 1} + \dots + a _ {0} = 0 \quad (a _ {1} \in A, a _ {0} \neq 0).
\]

这个方程表明 \(a_{0} \in By \subset b\)。将引理3应用于 M = B，可见 \(B/a_{0}B\) 是有限长度的 A-模；因此其商模 B/b 也具有有限长度。此外，B-模 b 含有有限生成的子模 \(a_{0}B\)；由于 \(b/a_{0}B\) 是有限长度的（作为 \(B/a_{0}B\) 的子模），因而是有限生成的，所以 b 确实是有限生成的 B-模。

上述结果首先表明 B 是诺特环。另一方面，若 p 是 B 的非零素理想，则 B/p 是整环且具有有限长度，因而是域（《代数》第VIII章，第6节，第4号，命题9），所以 p 是极大理想。

推论1. 对于 A 的每个素理想 p，B 中位于 p 之上的素理想集合是有限的。

先假设 \(\mathfrak{p} = (0)\)；则唯一的素理想 \(\mathfrak{q}\)（属于 \(\mathbf{B}\)）满足 \(\mathfrak{q} \cap \mathbf{A} = (0)\)，就是 (0)；否则，记 \(\mathbf{S} = \mathbf{A} - \{0\}\)，则 \(\mathbf{S}^{-1}\mathfrak{q}\) 是 \(\mathbf{S}^{-1}\mathbf{B}\) 的非零素理想（第II章，第2节，第5号，命题11），并且 \(\mathbf{S}^{-1}\mathbf{B}\) 正是 \(\mathbf{B}\) 的分式域，因为它是 \(\mathbf{L}\) 的包含 \(\mathbf{K}\) 的子环（《代数》第V章，第3节，第2号，命题3）；从而得到荒谬结论。若此时 \(\mathfrak{p} \neq (0)\)，则由命题5可知 \(\mathbf{B}/\mathfrak{p}\mathbf{B}\) 是在域 \(\mathbf{A}/\mathfrak{p}\) 上的有限维向量空间，因而是阿廷环，所以只有有限多个素理想（第IV章，第2节，第5号，命题9）；这说明 \(\mathbf{B}\) 中包含 \(\mathfrak{p}\) 的素理想只有有限多个。

推论2. A 在 L 中的整闭包是 Dedekind 整环。

这个整闭包是整闭的诺特整环，其所有非零素理想都是极大理想；因此只需应用第2节第1号定理。

特别地：

推论3. Dedekind 整环在其分式域的有限扩张中的整闭包是 Dedekind 整环。

命题6. 设 A 是 Dedekind 整环，K 为其分式域，L 是 K 的有限扩张，B 是 A 在 L 中的整闭包。设 p 是 A 的非零素理想，v 是 K 上相应的本性赋值，并且

\[
\mathrm{B} p = \prod_ {i} p _ {i} ^ {e (i)}
\]

理想 Bp 分解为素理想的乘积。则：

\begin{enumerate}
\def\labelenumi{(\alph{enumi})}
\item
  B 中位于 p 之上的素理想为 \(p_{i}\)，其中 \(e(i) > 0\)；
\item
  赋值 \(v_{i}\)（定义在 \(L\) 上，与这些理想 \(p_{i}\) 对应），在等价意义下，正是定义在 \(L\) 上并延拓 \(v\) 的赋值；
\item
  \([\mathbf{B} / \mathfrak{p}_i\colon \mathbf{A} / \mathfrak{p}] = f(v_i / v)\)；
\item
  \(e_i = e(v_i / v)\)（参见第VI章，第8节，第1号，定义1和2）。
\item
  说 B 的素理想 q 位于 p 之上，等价于 \(q \supset p\)，从而 \(q \supset Bp\)，并且 q 包含某个 \(p_i\)，满足 \(e(i) > 0\)（第II章，第1节，第1号，命题1）。
\item
  这由 (a) 以及第1节第8号命题12的推论得到。
\item
  \(v\) 的剩余域可与 \(A / p\) 等同，而 \(v_i\) 的剩余域可与 \(B / p_i\) 等同（第1节，第4号，命题6的推论1）。
\item
  设 \(a\)（相应地，\(a_i\)）分别是 \(v\)（相应地，\(v_i\)）的均匀化元。然后
\end{enumerate}

\[
\begin{array}{r l} a \mathrm{B} _ {p _ {i}} = a \mathrm{A} _ {p} \mathrm{B} _ {p _ {i}} = p \mathrm{A} _ {p} \mathrm{B} _ {p _ {i}} = p \mathrm{B}. \mathrm{B} _ {p _ {i}} & = \left(\prod_ {j} p _ {j} ^ {e (j)}\right) \mathrm{B} _ {p _ {i}} = \prod_ {j} \left(p _ {j} \mathrm{B} _ {p _ {i}}\right) ^ {e (j)} \\ & = \left(p _ {i} \mathrm{B} _ {p _ {i}}\right) ^ {e (i)} = a _ {i} ^ {e (i)} \mathrm{B} _ {p _ {i}} \end{array}
\]

由于 \(\mathfrak{p}_j\mathrm{B}_{\mathfrak{p}_i} = \mathrm{B}_{\mathfrak{p}_i}\) 对 \(j\neq i\) 成立，故由 (d)，\(e(v_{i} / v) = v_{i}(a)\)。

